% !TeX program = xelatex
\documentclass[UTF8,zihao=-4,fontset=fandol]{ctexart}
\newcommand{\ExamNumber}{1}
\usepackage{linear-algebra-exam}

\begin{document}
\ExamCover{9}{72}{60}

\Question{1}{单项选择}{2014 年数学一第 5 题、数学二第 7 题（均为原卷第 1 页）}{4}{***}
行列式
\[
\begin{vmatrix}
0&a&b&0\\
a&0&0&b\\
0&c&d&0\\
c&0&0&d
\end{vmatrix}
=\text{（\qquad）。}
\]
\Choices{$(ad-bc)^2$}{$-(ad-bc)^2$}{$a^2d^2-b^2c^2$}{$b^2c^2-a^2d^2$}

\Solution
\Known 求含 $a,b,c,d$ 的四阶行列式。每行有两个零，其中第 $1$、$3$ 行的零位相同，第 $2$、$4$ 行的零位也相同。
\Method 把非零元素集中到两个对角块，四阶行列式就可以用两个二阶行列式相乘来计算。每次单独交换两行或两列，整体的行列式都变号。
\Derive 先交换第 $1$、$2$ 列，行列式变号。再交换第 $2$、$3$ 行，同时交换第 $2$、$3$ 列，这两次交换的符号相互抵消。所得行列式为
\[
\begin{vmatrix}
a&b&0&0\\
c&d&0&0\\
0&0&a&b\\
0&0&c&d
\end{vmatrix}
=\begin{vmatrix}a&b\\c&d\end{vmatrix}^{\!2}
=(ad-bc)^2.
\]
总共交换三次，所以原行列式为 $-(ad-bc)^2$。
第一次交换后，四行依次为 $(a,0,b,0)$、$(0,a,0,b)$、$(c,0,d,0)$、$(0,c,0,d)$。第 $1$、$3$ 行的非零位置恰好是第 $1$、$3$ 列，所以把这两行和这两列同时排到最前面，就分出了两个相同的二阶块。分块对角行列式取乘积，是因为前两行只能从前两列取项，后两行只能从后两列取项。

另外三个选项可用同一组数排除。取 $a=d=2,\ b=c=1$，四个选项依次给出 $9,-9,15,-15$。按上面的交换计算，原行列式为 $-9$。
\Verify 选 \textbf{B}。取 $a=d=1,\ b=c=0$，原矩阵是单位矩阵交换前两行，行列式为 $-1$，可核对负号。

\Question{2}{单项选择}{2023 年数学二第 8 题（原卷第 2 页）}{4}{****}
设 $A,B$ 为 $n$ 阶可逆矩阵，$E$ 为 $n$ 阶单位矩阵，$M^*$ 为矩阵 $M$ 的伴随矩阵，则 $\begin{pmatrix}A&E\\O&B\end{pmatrix}^{*}=$（\quad）。
\Choices
  {$\begin{pmatrix}|A|B^*&-B^*A^*\\O&|B|A^*\end{pmatrix}$}
  {$\begin{pmatrix}|A|B^*&-A^*B^*\\O&|B|A^*\end{pmatrix}$}
  {$\begin{pmatrix}|B|A^*&-B^*A^*\\O&|A|B^*\end{pmatrix}$}
  {$\begin{pmatrix}|B|A^*&-A^*B^*\\O&|A|B^*\end{pmatrix}$}

\Solution
\Known $M$ 是分块上三角矩阵，对角块 $A,B$ 的行列式均非零。要求 $M$ 的伴随矩阵。
\Method 由 $M^*=|M|M^{-1}$，可以把问题拆成求分块行列式和求逆矩阵。逆矩阵右上角应消去乘积中由 $E$ 引入的那一项，因子顺序由分块乘法决定。
\Derive $|M|=|A|\,|B|$。计算
\[
\begin{pmatrix}A&E\\O&B\end{pmatrix}
\begin{pmatrix}A^{-1}&-A^{-1}B^{-1}\\O&B^{-1}\end{pmatrix}
=\begin{pmatrix}E&O\\O&E\end{pmatrix}.
\]
右上块之和是 $-AA^{-1}B^{-1}+B^{-1}=O$。因此第二个矩阵是 $M^{-1}$。分别代入 $A^{-1}=\frac{A^*}{|A|}$ 和 $B^{-1}=\frac{B^*}{|B|}$，并乘上 $|A|\,|B|$，得到
\[
M^*=\begin{pmatrix}|B|A^*&-A^*B^*\\O&|A|B^*\end{pmatrix}.
\]
右上角的次序可以从未知量直接推出来。令逆矩阵右上角为 $Y$，右下角为 $B^{-1}$，则右上角乘积要求 $AY+EB^{-1}=O$。左乘 $A^{-1}$ 才得到 $Y=-A^{-1}B^{-1}$。例如 $A=\operatorname{diag}(1,2)$、$B=\begin{pmatrix}1&1\\0&1\end{pmatrix}$ 时，$A^*B^*$ 的右上元为 $-2$，而 $B^*A^*$ 的右上元为 $-1$，所以两个因子一般不能互换。
\Verify 选 \textbf{D}。用 $M$ 左乘此结果，右上块为 $-AA^*B^*+|A|B^*=O$，其余各块与 $|M|E_{2n}$ 一致。

\Question{3}{单项选择}{2026 年数学二第 9 题（原卷第 2 页）}{4}{**}
设矩阵
\[
A=\begin{pmatrix}1&0&1\\0&0&1\\1&1&3\\1&1&1\end{pmatrix},
\qquad C=\begin{pmatrix}2&0\\1&1\\1&1\\a&b\end{pmatrix}\text{。}
\]
若存在矩阵 $B$ 满足 $AB=C$，则（\quad）。
\Choices
  {$a=-1,\ b=-1$}
  {$a=2,\ b=2$}
  {$a=-1,\ b=2$}
  {$a=2,\ b=-1$}

\Solution
\Known 已知 $A$ 及乘积 $C=AB$ 的前三行。要求 $C$ 第四行的两个元素 $a,b$。
\Method 乘积矩阵的第 $i$ 行等于 $A$ 的第 $i$ 行右乘 $B$。若 $A$ 的最后一行能由前面的行相加减得到，右乘 $B$ 后，同样的等式也成立。
\Derive $A$ 的第三行减去第二行的 $2$ 倍为
\[
(1,1,3)-2(0,0,1)=(1,1,1)\text{，}
\]
恰好是第四行。右乘 $B$ 并利用分配律，$C$ 的第四行亦为其第三行减去第二行的 $2$ 倍，所以
\[
(a,b)=(1,1)-2(1,1)=(-1,-1)\text{。}
\]
把存在性检查展开，可设 $B$ 的三行为 $r_1,r_2,r_3$。由 $A$ 的第二行可知 $r_3=(1,1)$。第一行给出 $r_1+r_3=(2,0)$，所以 $r_1=(1,-1)$。第三行给出 $r_1+r_2+3r_3=(1,1)$，解得 $r_2=(-3,-1)$。最后由第四行计算得到 $r_1+r_2+r_3=(-1,-1)$。
\Verify 选 \textbf{A}。取 $B=\begin{pmatrix}1&-1\\-3&-1\\1&1\end{pmatrix}$，直接相乘即得含 $a=b=-1$ 的 $C$，因此题设的存在条件也得到验证。

\Question{4}{填空}{2012 年数学二第 14 题（原卷第 2 页）}{8}{**}
设 $A$ 为 $3$ 阶矩阵，$|A|=3$，$A^*$ 为 $A$ 的伴随矩阵，若交换 $A$ 的第 $1$ 行与第 $2$ 行得矩阵 $B$，则 $|BA^*|=\blank$。

\Solution
\Known 已知 $A$ 的行列式及 $B$ 的构造方式，要求乘积矩阵 $BA^*$ 的行列式。
\Method 交换两行使行列式变号，先求 $|B|$。已知 $AA^*=|A|E$，所以对这个等式取行列式即可求出 $|A^*|$。
\Derive 交换两行后 $|B|=-|A|=-3$。在 $AA^*=3E$ 两边取行列式，利用乘积的行列式等于行列式的乘积，得
\[
3|A^*|=|3E|=3^3,\qquad |A^*|=9.
\]
$3E$ 是三阶对角矩阵，所以要乘三次 $3$。于是 $|BA^*|=|B|\,|A^*|=-27$。
伴随矩阵的恒等式来自按行展开。$AA^*$ 的第 $(i,j)$ 元是 $A$ 的第 $i$ 行元素与第 $j$ 行代数余子式的乘积之和。若 $i=j$，这个和就是 $|A|$；若 $i\ne j$，这个和可看作把第 $j$ 行换成第 $i$ 行后按第 $j$ 行展开，两行相同，行列式为零。故 $AA^*=|A|E$，这里可以用它计算 $|A^*|$。
\Verify 填 $\boxed{-27}$。也可直接检查 $BA^*$。它由 $AA^*=3E$ 交换前两行得到，因此行列式等于 $-|3E|=-27$。

\Question{5}{填空}{2016 年数学一第 13 题（原卷第 2 页）}{8}{***}
行列式
\[
\begin{vmatrix}
\lambda&-1&0&0\\
0&\lambda&-1&0\\
0&0&\lambda&-1\\
4&3&2&\lambda+1
\end{vmatrix}
=\blank
\]

\Solution
\Known 求含参数 $\lambda$ 的四阶行列式。前三行除主对角线及其上方的 $-1$ 外均为零。
\Method 先按第一列展开。这一列只有首尾两项，其中 $4$ 的余子式是三角行列式，另一项仍有相同的稀疏结构，可再按第一列展开。
\Derive 记原行列式为 $D$，则
\[
\begin{aligned}
D&=\lambda\begin{vmatrix}\lambda&-1&0\\0&\lambda&-1\\3&2&\lambda+1\end{vmatrix}
 -4\begin{vmatrix}-1&0&0\\\lambda&-1&0\\0&\lambda&-1\end{vmatrix}\\
 &=\lambda\left[\lambda
 \begin{vmatrix}\lambda&-1\\2&\lambda+1\end{vmatrix}+3\right]+4\\
 &=\lambda\bigl[\lambda(\lambda^2+\lambda+2)+3\bigr]+4.
\end{aligned}
\]
第一步中 $4$ 的代数余子式符号为 $(-1)^{4+1}=-1$，它的余子式值为 $(-1)^3=-1$，所以贡献 $+4$。第二步中的 $3$ 对应三阶行列式的位置 $(3,1)$，符号为正，余子式的值为 $1$。
把内层二阶行列式展开为 $\lambda(\lambda+1)-(-1)\cdot2=\lambda^2+\lambda+2$。它在公式中先被内层展开乘一次 $\lambda$，再被外层展开乘一次 $\lambda$，所以这一部分是 $\lambda^2(\lambda^2+\lambda+2)$；第二层的 $3$ 只乘外层的 $\lambda$。合并时可分别检查这两部分与常数 $4$，避免把 $3\lambda$ 写成 $3\lambda^2$。
\Verify 填 $\boxed{\lambda^4+\lambda^3+2\lambda^2+3\lambda+4}$。令 $\lambda=0$，原式按首列展开只剩 $4$ 的贡献，和常数项一致。

\Question{6}{填空}{2019 年数学二第 14 题（原卷第 2 页）}{8}{**}
已知矩阵 $A=\begin{pmatrix}1&-1&0&0\\-2&1&-1&1\\3&-2&2&-1\\0&0&3&4\end{pmatrix}$，$A_{ij}$ 表示 $|A|$ 中 $(i,j)$ 元的代数余子式，则 $A_{11}-A_{12}=\blank$。

\Solution
\Known 求第一行前两个代数余子式之差。代数余子式是删去对应行列后所得的行列式，再乘位置符号 $(-1)^{i+j}$。
\Method $(1,1)$ 的符号为正，$(1,2)$ 的符号为负，因此相减会变成两个三阶行列式相加。它们的后两列相同，可以先合并第一列。
\Derive 删去第一行，并分别删去第一、第二列，得
\[
A_{11}-A_{12}=
\begin{vmatrix}1&-1&1\\-2&2&-1\\0&3&4\end{vmatrix}
+\begin{vmatrix}-2&-1&1\\3&2&-1\\0&3&4\end{vmatrix}
=\begin{vmatrix}-1&-1&1\\1&2&-1\\0&3&4\end{vmatrix}.
\]
沿第一列展开，第一项为 $-1(8+3)=-11$，第二项为 $-1(-4-3)=7$，两项相加得 $-4$。
若分别计算，则第一个三阶行列式沿第一列展开得 $(8+3)+2(-4-3)=-3$；第二个三阶行列式得 $-2(8+3)-3(-4-3)=-1$。因此 $A_{11}=-3$、$A_{12}=1$。第二个行列式是余子式，必须再乘 $(-1)^{1+2}$ 才得到代数余子式。
\Verify 填 $\boxed{-4}$。也可以按原矩阵第一行展开。该行只有 $1,-1$ 两项非零，所以这个余子式之差正好是 $|A|$。

\Question{7}{填空}{2022 年数学一第 15 题（原卷第 2 页）}{8}{***}
已知矩阵 $A$ 和 $E-A$ 可逆，其中 $E$ 为单位矩阵。若矩阵 $B$ 满足 $[E-(E-A)^{-1}]B=A$，则 $B-A=\blank$。

\Solution
\Known $A$ 与 $E-A$ 均可逆，给定一个含 $B$ 的矩阵等式，要求 $B-A$。
\Method 方括号内有 $(E-A)^{-1}$。在等式两边同时左乘 $E-A$，可以用相邻的两个互逆矩阵消去这一项。
\Derive 左乘 $E-A$ 后，左侧的系数为
\[
(E-A)[E-(E-A)^{-1}]=(E-A)-E=-A,
\]
所以 $-AB=(E-A)A$。右侧按分配律展开为 $A-A^2=A(E-A)$，故 $-AB=A(E-A)$。因为 $A$ 可逆，在两边左乘 $A^{-1}$ 得 $-B=E-A$，即 $B=A-E$。
也可直接因式分解方括号。由 $E=(E-A)^{-1}(E-A)$，
\[
E-(E-A)^{-1}=-(E-A)^{-1}A.
\]
右端由两个可逆矩阵相乘而成，说明原方程的左系数可逆，所得 $B$ 唯一。原解法里 $A$ 与 $E-A$ 可以换序，是因为两种乘法展开都等于 $A-A^2$。这一点不能推广到任意两个矩阵。
\Verify 填 $\boxed{-E}$。代入 $B=-(E-A)$，左侧化为 $-(E-A)+E=A$，满足原等式。

\Question{8}{填空}{1994 年数学一（见 2013 年数学二配套解析 PDF 第 26 页）}{8}{**}
已知 $\alpha=(1,2,3)$，$\beta=\left(1,\frac{1}{2},\frac{1}{3}\right)$，设 $A=\alpha^{\mathrm T}\beta$，其中 $\alpha^{\mathrm T}$ 是 $\alpha$ 的转置，则 $A^n=\blank$。

\Solution
\Known $\alpha,\beta$ 是行矩阵，$A$ 是列矩阵 $\alpha^{\mathrm T}$ 与行矩阵 $\beta$ 相乘得到的三阶矩阵。要求正整数 $n$ 次幂。
\Method 相乘两个 $A$ 时，中间碰到的是行乘列 $\beta\alpha^{\mathrm T}$。这是一个数，可先把它算出来，之后每多乘一个 $A$ 都会产生同样的因子。
\Derive $\beta\alpha^{\mathrm T}=1+1+1=3$。由结合律，$A^2=\alpha^{\mathrm T}(\beta\alpha^{\mathrm T})\beta=3A$。若 $A^n=3^{n-1}A$，则 $A^{n+1}=3^{n-1}A^2=3^nA$，所以
\[
A^n=3^{n-1}\begin{pmatrix}
1&\frac{1}{2}&\frac{1}{3}\\
2&1&\frac{2}{3}\\
3&\frac{3}{2}&1
\end{pmatrix}.
\]
尺寸也提供了入手线索。$\alpha^{\mathrm T}\beta$ 是 $3\times3$ 矩阵，它的第 $i$ 行是 $\alpha$ 的第 $i$ 个分量乘 $\beta$。反过来 $\beta\alpha^{\mathrm T}$ 是 $1\times1$，可作为标量提出。归纳的起点是 $A^1=3^0A$，每向下一次幂都只有中间的行乘列增加一个因子 $3$。
\Verify 取 $n=1$ 回到列矩阵乘行矩阵的原式，取 $n=2$ 对应上面直接算出的 $3A$。

\Question{9}{解答}{2013 年数学一第 20 题（第 3 页）、数学二第 22 题（第 4 页）}{20}{****}
设 $A=\begin{pmatrix}1&a\\1&0\end{pmatrix}$，$B=\begin{pmatrix}0&1\\1&b\end{pmatrix}$。当 $a,b$ 为何值时，存在矩阵 $C$ 使得 $AC-CA=B$，并求所有矩阵 $C$。

\Solution
\Known $A,B$ 是二阶矩阵，题设的两个乘法都要求 $C$ 也是二阶矩阵。先找能使等式成立的 $a,b$，再求此时全部的 $C$。
\Method 令 $C=\begin{pmatrix}x&y\\z&w\end{pmatrix}$，直接做矩阵乘法。两个矩阵相等意味着对应位置的元素相等，可以逐项代入，判断哪些未知数被约束、哪些可以任取。
\Derive 按矩阵乘法计算，
\[
AC-CA=\begin{pmatrix}
az-y & y+a(w-x)\\
x-z-w & y-az
\end{pmatrix}\text{。}
\]
左上角等于 $0$，所以 $y=az$。右下角恰好是左上角的相反数，因此 $b=0$。左下角等于 $1$，所以 $x-w=1+z$。把这两个等式代入右上角，得到
\[
1=y+a(w-x)=az-a(1+z)=-a\text{，}
\]
因此 $a=-1$。此时 $y=-z$、$x=w+z+1$，其余等式对 $z,w$ 没有额外限制。
四个位置的比较可写为
\[
az-y=0,\qquad y+a(w-x)=1,\qquad x-z-w=1,\qquad y-az=b.
\]
第一式和第四式相加当即给出 $b=0$。参数确定后，第一式和第三式分别固定了 $y$ 与 $x-w$，其余两式由这两个关系自动满足。这解释了为什么 $z,w$ 可独立取值，也解释了最后写出的解能覆盖所有可能的 $C$。
\Verify 当 $a=-1,\ b=0$ 时，所有解为
\[
C=\begin{pmatrix}s+t+1&-t\\t&s\end{pmatrix},\qquad s,t\in\mathbb R\text{。}
\]
将它代回差矩阵，任意 $s,t$ 都给出 $\begin{pmatrix}0&1\\1&0\end{pmatrix}$，因此上述条件既必要也足够。
\end{document}
